\SourceLeaf{212}lyse, durch welche einige Geometer auf der Oberfläche der Kugel den Werth $-2\pi k$, oder den Mittelwerth zwischen dem innen und aussen geltenden, herausgebracht haben, kann ich, insofern der Begriff von Differentialquotienten in seiner mathematischen Reinheit aufgefasst wird, eine Zulässigkeit nicht einräumen.

\Article{9}
Das im vorhergehenden Beispiel gefundene Resultat ist nur ein einzelner Fall des allgemeinen Theorems, nach welchem, wenn der Punkt $O$ sich im Innern der wirkenden Masse befindet, der Werth von $\frac{\mathrm{d}\mathrm{d}V}{\mathrm{d}x^2}+\frac{\mathrm{d}\mathrm{d}V}{\mathrm{d}y^2}+\frac{\mathrm{d}\mathrm{d}V}{\mathrm{d}z^2}$ äqual wird dem Producte aus $-4\pi$ in die in $O$ Statt findende Dichtigkeit. Die befriedigendste Art, diesen wichtigen Lehrsatz zu begründen, scheint folgende zu sein.

Wir nehmen an, dass die Dichtigkeit $k$ sich innerhalb $t$ nirgends sprungs\-weise ändere, oder dass sie eine mit $f(a,b,c)$ zu bezeichnende Function von $a$, $b$, $c$ sei, deren Werth sich innerhalb $t$ überall nach der Stetigkeit ändert, ausserhalb $t$ hingegen $=0$ wird.

Es sei $t'$ der Raum, in welchen $t$ übergeht, wenn die erste Coordinate jedes Punktes der Grenzfläche um die Grösse $e$ vermindert, oder was dasselbe ist, wenn die Grenzfläche parallel mit der ersten Coordinatenaxe um $e$ rückwärts bewegt wird; es bestehe $t$ aus den Räumen $t^0$ und $\theta$, $t'$ aus $t^0$ und $\theta'$, so dass $t^0$ der ganze Raum ist, welcher $t$ und $t'$ gemeinschaftlich bleibt. Wir betrachten die drei Integrale
\[
\begin{array}{ll}
\displaystyle\int\frac{f(a,b,c)(a-x)\mathrm{d}t}{((a-x)^2+(b-y)^2+(c-z)^2)^{\frac{3}{2}}}&.\ .\ .\ .\ .\ .\ .\ .\quad(1)\\[8pt]
\displaystyle\int\frac{f(a,b,c)(a-x-e)\mathrm{d}t}{((a-x-e)^2+(b-y)^2+(c-z)^2)^{\frac{3}{2}}}&.\ .\ .\ .\ .\ .\ .\quad(2)\\[8pt]
\displaystyle\int\frac{f(a+e,b,c)(a-x)\mathrm{d}t}{((a-x)^2+(b-y)^2+(c-z)^2)^{\frac{3}{2}}}&.\ .\ .\ .\ .\ .\ .\ .\quad(3)
\end{array}
\]
wo das Integral (1), über den ganzen Raum $t$ ausgedehnt, der Werth von $\frac{\mathrm{d}V}{\mathrm{d}x}$ oder $X$ in dem Punkte $O$ sein wird. Das Integral (2), gleichfalls über ganz $t$ ausgedehnt, wird der Werth von $\frac{\mathrm{d}V}{\mathrm{d}x}$ in demjenigen Punkte sein, dessen Coordinaten $x+e$, $y$, $z$ sind, welchen Werth wir mit $X+\xi$ bezeichnen wollen. Offenbar ist mit diesem Integrale ganz identisch das Integral (3), über den ganzen Raum $t'$ ausgedehnt. Ist also
\SourceLeaf{213}%
\[
\begin{array}{rlcl}
\text{das Integral (1), ausgedehnt über}&t^0&.\ .\ .\ .\ .\ .\ .\ .&l\\
\text{über}&\theta&.\ .\ .\ .\ .\ .\ .\ .&\lambda\\
\text{das Integral (3), ausgedehnt über}&t^0&.\ .\ .\ .\ .\ .\ .\ .&l'\\
\text{über}&\theta'&.\ .\ .\ .\ .\ .\ .\ .&\lambda'
\end{array}
\]
so wird $X=l+\lambda$, $X+\xi=l'+\lambda'$.

Setzen wir $f(a+e,b,c)-f(a,b,c)=\Delta k$, so ist das Integral
\[
\int\frac{\frac{\Delta k}{e}(a-x)\mathrm{d}t}{((a-x)^2+(b-y)^2+(c-z)^2)^{\frac{3}{2}}}\quad .\ .\ .\ .\ .\ .\ .\ .\quad(4)
\]
über $t^0$ ausgedehnt, $=\frac{l'-l}{e}$.

Die bisherigen Resultate gelten allgemein für jede Lage von $O$: bei der weitern Entwickelung soll der Fall, wo $O$ in der Oberfläche selbst liegt, ausgeschlossen sein, oder angenommen werden, dass $O$ in messbarer Entfernung von der Oberfläche, innerhalb oder ausserhalb $t$ liege.

Lassen wir nun $e$ unendlich klein werden, so sind die Räume $\theta$, $\theta'$ zwei unendlich schmale an der Oberfläche von $t$ anliegende Raumschichten; zerlegen wir diese Oberfläche in Elemente $\mathrm{d}s$, und bezeichnen mit $\alpha$ den Winkel, welchen eine in $\mathrm{d}s$ nach aussen errichtete Normale mit der ersten Coordinatenaxe macht, so wird $\alpha$ offenbar spitz sein überall, wo die Oberfläche von $t$ an $\theta$ grenzt, stumpf hingegen da, wo sie an $\theta'$ grenzt. Die Elemente von $\theta$ werden also ausgedrückt werden durch $e\cos\alpha\mathrm{d}s$, die Elemente von $\theta'$ hingegen durch $-e\cos\alpha\mathrm{d}s$, woraus man leicht schliesst, dass $\frac{\lambda-\lambda'}{e}$ übergeht in das Integral
\[
\int\frac{f(a,b,c)(a-x)\cos\alpha.\mathrm{d}s}{((a-x)^2+(b-y)^2+(c-z)^2)^{\frac{3}{2}}}
\]
oder was dasselbe ist, in dieses
\[
\int\frac{k(a-x)\cos\alpha.\mathrm{d}s}{r^3}
\]
durch die ganze Oberfläche ausgedehnt, wo unter $k$ die an dem Elemente $\mathrm{d}s$ Statt findende Dichtigkeit zu verstehen ist.

Unter Voraussetzung eines unendlich kleinen Werthes von $e$ wird ferner $\frac{\Delta k}{e}$ übergehen in den Werth des partiellen Differentialquotienten $\frac{\mathrm{d}f(a,b,c)}{\mathrm{d}a}$ oder $\frac{\mathrm{d}k}{\mathrm{d}a}$, und der Werth des Integrals (4) oder $\frac{(l'-l)}{e}$ in das Integral
\SourceLeaf{214}%
\[
\int\frac{\frac{\mathrm{d}k}{\mathrm{d}a}.(a-x)\mathrm{d}t}{r^3}
\]
durch den ganzen Raum $t$ ausgedehnt.

Endlich ist, für ein unendlich kleines $e$, $\frac{l'-l}{e}-\frac{\lambda-\lambda'}{e}$ oder $\frac{\xi}{e}$, nichts anderes, als der Werth des partiellen Differentialquotienten $\frac{\mathrm{d}X}{\mathrm{d}x}$ oder $\frac{\mathrm{d}\mathrm{d}V}{\mathrm{d}x^2}$. Wir haben folglich das einfache Resultat
\[
\frac{\mathrm{d}\mathrm{d}V}{\mathrm{d}x^2}=\frac{\mathrm{d}X}{\mathrm{d}x}=\int\frac{\frac{\mathrm{d}k}{\mathrm{d}a}.(a-x)\mathrm{d}t}{r^3}-\int\frac{k(a-x)\cos\alpha.\mathrm{d}s}{r^3}
\]
wo die erste Integration über den ganzen Raum $t$, die zweite über die ganze Oberfläche desselben auszudehnen ist.

Dieses Resultat ist gültig, wie nahe auch $O$ der Oberfläche auf der innern oder äussern Seite liegen mag, nur nicht in der Oberfläche selbst, wo vielmehr $\frac{\mathrm{d}X}{\mathrm{d}x}$ zwei verschiedene Werthe haben wird. Das erste Integral ändert sich zwar beim Durchgange durch die Oberfläche nach der Stetigkeit, hingegen ändert sich $-\int\frac{k(a-x)\cos\alpha\mathrm{d}s}{r^3}$ nach einem weiter unten zu beweisenden Theorem beim Übergange von einem innern der Oberfläche unendlich nahen Punkte nach einem äussern um die endliche Grösse $4\pi k\cos\alpha$, wo $k$ und $\alpha$ sich auf die Durchgangsstelle beziehen, und eben so gross wird der Unterschied der beiden daselbst Statt findenden Werthe von $\frac{\mathrm{d}X}{\mathrm{d}x}$ sein.

\Article{10}
Auf ähnliche Weise wird, wenn $\beta$ und $\gamma$ in Beziehung auf die zweite und dritte Coordinatenaxe dieselbe Bedeutung haben, wie $\alpha$ in Beziehung auf die erste, und für die Lage von $O$ dieselbe Beschränkung gilt, wie vorhin,
\[
\begin{aligned}
\frac{\mathrm{d}Y}{\mathrm{d}y}&=\int\frac{\frac{\mathrm{d}k}{\mathrm{d}b}(b-y)\mathrm{d}t}{r^3}-\int\frac{k(b-y)\cos\beta.\mathrm{d}s}{r^3}\\
\frac{\mathrm{d}Z}{\mathrm{d}z}&=\int\frac{\frac{\mathrm{d}k}{\mathrm{d}c}(c-z)\mathrm{d}t}{r^3}-\int\frac{k(c-z)\cos\gamma.\mathrm{d}s}{r^3}
\end{aligned}
\]
Erwägen wir nun, dass
\[
\frac{\mathrm{d}k}{\mathrm{d}a}.\frac{a-x}{r}+\frac{\mathrm{d}k}{\mathrm{d}b}.\frac{b-y}{r}+\frac{\mathrm{d}k}{\mathrm{d}c}.\frac{c-z}{r}
\]
\SourceLeaf{215}%
nichts anderes ist, als der Werth des Differentialquotienten $\frac{\mathrm{d}k}{\mathrm{d}r}$, insofern in dieser Differentiation nur die Länge von $r$ als veränderlich, die Richtung aber als constant betrachtet wird; ferner, dass
\[
\frac{a-x}{r}.\cos\alpha+\frac{b-y}{r}.\cos\beta+\frac{c-z}{r}.\cos\gamma=\cos\psi
\]
wird, wenn $\psi$ den Winkel bezeichnet, welchen die nach aussen gerichtete Normale in $\mathrm{d}s$ mit der verlängerten geraden Linie $r$ macht, so erhellt, dass, wenn das Integral
\[
\int\frac{\frac{\mathrm{d}k}{\mathrm{d}r}}{rr}.\mathrm{d}t
\]
über den ganzen Raum $t$ erstreckt mit $M$, das Integral
\[
\int\frac{k\cos\psi}{rr}.\mathrm{d}s
\]
durch die ganze Oberfläche von $t$ ausgedehnt mit $N$ bezeichnet wird,
\[
\frac{\mathrm{d}\mathrm{d}V}{\mathrm{d}x^2}+\frac{\mathrm{d}\mathrm{d}V}{\mathrm{d}y^2}+\frac{\mathrm{d}\mathrm{d}V}{\mathrm{d}z^2}=M-N
\]
sein wird.

Um die erstere Integration auszuführen, beschreiben wir um den Mittelpunkt $O$ mit dem Halbmesser $1$ eine Kugelfläche, und zerlegen dieselbe in Elemente $\mathrm{d}\sigma$. Die von $O$ durch alle Punkte der Peripherie von $\mathrm{d}\sigma$ geführten und unbestimmt verlängerten geraden Linien bilden eine Kegelfläche (im weitern Sinne des Worts), wodurch aus dem ganzen $t$ ein Raum (nach Umständen aus mehrern getrennten Stücken bestehend) ausgeschieden wird, und wovon $rr\mathrm{d}\sigma.\mathrm{d}r$ ein unbestimmtes Element ist. Derjenige Theil von $M$, welcher sich auf diesen Raum bezieht, wird folglich durch $\mathrm{d}\sigma.\int\frac{\mathrm{d}k}{\mathrm{d}r}.\mathrm{d}r$ ausgedrückt werden, wenn diese Integration durch alle in $t$ fallenden Theile einer durch $O$ und einen Punkt von $\mathrm{d}\sigma$ gehenden soweit als nöthig verlängerten geraden Linie $r$ erstreckt wird. Nehmen wir nun an, diese gerade Linie schneide die Oberfläche von $t$ der Reihe nach in $O'$, $O''$, $O'''$, $O''''$ u. s. f.; bezeichnen mit $r'$, $r''$, $r'''$, $r''''$ u. s. f. die Werthe von $r$ in diesen Punkten; mit $\mathrm{d}s'$, $\mathrm{d}s''$, $\mathrm{d}s'''$, $\mathrm{d}s''''$ u. s. f. die entsprechenden durch den Elementarkegel aus der Oberfläche von $t$ ausgeschiedenen Elemente; mit $k'$, $k''$, $k'''$, $k''''$ u. s. f. die Werthe von $k$, und mit $\psi'$, $\psi''$, $\psi'''$, $\psi''''$ u. s. f. die Werthe von $\psi$ an diesen Elementen: so übersieht man leicht, dass
\SourceLeaf{216}%

I. für den Fall, wo $O$ innerhalb $t$ liegt, die Anzahl jener Punkte ungerade, und die Integration $\int\frac{\mathrm{d}k}{\mathrm{d}r}.\mathrm{d}r$ von $r=0$ bis $r=r'$, dann von $r=r''$ bis $r=r'''$ u. s. f. auszuführen sein wird, woraus also, wenn die Dichtigkeit in $O$ mit $k^0$ bezeichnet wird, hervorgeht
\[
\int\frac{\mathrm{d}k}{\mathrm{d}r}.\mathrm{d}r=-k^0+k'-k''+k'''-k''''+\text{ u. s. f.}
\]
Da die Winkel $\psi'$, $\psi''$, $\psi'''$, $\psi''''$ u. s. f. offenbar abwechselnd spitz und stumpf sind, so wird
\[
\begin{aligned}
\mathrm{d}s'.\cos\psi'&=+r'r'\mathrm{d}\sigma\\
\mathrm{d}s''.\cos\psi''&=-r''r''\mathrm{d}\sigma\\
\mathrm{d}s'''.\cos\psi'''&=+r'''r'''\mathrm{d}\sigma\\
\mathrm{d}s''''.\cos\psi''''&=-r''''r''''\mathrm{d}\sigma
\end{aligned}
\]
u. s. f. und folglich
\[
\begin{aligned}
\mathrm{d}\sigma\int\frac{\mathrm{d}k}{\mathrm{d}r}.\mathrm{d}r&=-k^0\mathrm{d}\sigma+\frac{k'\cos\psi'}{r'r'}\mathrm{d}s'+\frac{k''\cos\psi''}{r''r''}\mathrm{d}s''+\frac{k'''\cos\psi'''}{r'''r'''}+\text{ u. s. f.}\\
&=-k^0\mathrm{d}\sigma+\sum\frac{k\cos\psi}{rr}\mathrm{d}s
\end{aligned}
\]
indem die Summation auf alle $\mathrm{d}s$ ausgedehnt wird, welche dem Element $\mathrm{d}\sigma$ entsprechen. Durch Integration über sämmtliche $\mathrm{d}\sigma$ erhält man also
\[
M=-4\pi k^0+\int\frac{k\cos\psi}{rr}\mathrm{d}s
\]
wo das Integral über die ganze Oberfläche erstreckt werden muss, oder $M=-4\pi k^0+N$. Es wird folglich
\[
\frac{\mathrm{d}\mathrm{d}V}{\mathrm{d}x^2}+\frac{\mathrm{d}\mathrm{d}V}{\mathrm{d}y^2}+\frac{\mathrm{d}\mathrm{d}V}{\mathrm{d}z^2}=-4\pi k^0
\]

II. Für den Fall, wo $O$ ausserhalb $t$ liegt, hat man nur diejenigen $\mathrm{d}\sigma$ in Betracht zu ziehen, für welche die durch $O$ und einen Punkt von $\mathrm{d}\sigma$ gelegte gerade Linie den Raum $t$ wirklich trifft; die Anzahl der Punkte $O'$, $O''$, $O'''$ u. s. f. wird hier immer gerade sein, und die Winkel $\psi'$, $\psi''$, $\psi'''$ u. s. f. abwechselnd stumpf und spitz, also
\[
\begin{aligned}
\mathrm{d}s'.\cos\psi'&=-r'r'\mathrm{d}\sigma,\\
\mathrm{d}s''.\cos\psi''&=+r''r''\mathrm{d}\sigma,\\
\mathrm{d}s'''\cos\psi'''&=-r'''r'''\mathrm{d}\sigma\text{ u. s. f.}
\end{aligned}
\]
Da nun hier die Integration $\int\frac{\mathrm{d}k}{\mathrm{d}r}.\mathrm{d}r$ von $r=r'$ bis $r=r''$, dann von $r=r'''$
\SourceLeaf{217}%
bis $r=r''''$ u. s. f. ausgeführt werden muss, so ergibt sich
\[
\begin{aligned}
\mathrm{d}\sigma.\int\frac{\mathrm{d}k}{\mathrm{d}r}.\mathrm{d}r&=\frac{k'\cos\psi'}{r'r'}.\mathrm{d}s'+\frac{k''\cos\psi''}{r''r''}.\mathrm{d}s''\\
&\quad{}+\frac{k'''\cos\psi'''}{r'''r'''}.\mathrm{d}s'''+\text{u. s. f.}\\
&=\sum\frac{k\cos\psi}{rr}\mathrm{d}s
\end{aligned}
\]
und nach der zweiten Integration durch alle in Betracht kommenden $\mathrm{d}\sigma$
\[
M=\int\frac{k\cos\psi}{rr}\mathrm{d}s=N
\]
folglich, wie ohnehin bekannt ist
\[
\frac{\mathrm{d}\mathrm{d}V}{\mathrm{d}x^2}+\frac{\mathrm{d}\mathrm{d}V}{\mathrm{d}y^2}+\frac{\mathrm{d}\mathrm{d}V}{\mathrm{d}z^2}=0
\]

\Article{11}
Obgleich in unsrer Beweisführung angenommen ist, dass die Dichtigkeit sich in dem \textit{ganzen} Raum $t$ nach der Stetigkeit ändere, so ist doch zur Gültigkeit unsers Resultats diese Bedingung nicht nothwendig, sondern es wird bloss erfordert, dass in dem Punkte $O$ die Dichtigkeit nach allen Seiten zu nach der Stetigkeit sich ändere, oder dass $O$ innerhalb eines wenn auch noch so kleinen dieser Bedingung Genüge leistenden Raumes liege. Setzen wir nemlich das Potential der in \textit{diesem} Raume enthaltenen Masse $=V'$, das Potential der übrigen ausserhalb desselben befindlichen Massen $=V''$, so wird das ganze Potential $V=V'+V''$, und da nach dem vorhergehenden Artikel
\[
\begin{aligned}
\frac{\mathrm{d}\mathrm{d}V'}{\mathrm{d}x^2}+\frac{\mathrm{d}\mathrm{d}V'}{\mathrm{d}y^2}+\frac{\mathrm{d}\mathrm{d}V'}{\mathrm{d}z^2}&=-4\pi k^0\\
\frac{\mathrm{d}\mathrm{d}V''}{\mathrm{d}x^2}+\frac{\mathrm{d}\mathrm{d}V''}{\mathrm{d}y^2}+\frac{\mathrm{d}\mathrm{d}V''}{\mathrm{d}z^2}&=0
\end{aligned}
\]
ist, so wird
\[
\frac{\mathrm{d}\mathrm{d}V}{\mathrm{d}x^2}+\frac{\mathrm{d}\mathrm{d}V}{\mathrm{d}y^2}+\frac{\mathrm{d}\mathrm{d}V}{\mathrm{d}z^2}=-4\pi k^0
\]
Fehlt hingegen diese Bedingung in dem Punkte $O$, und liegt also dieser in der Scheidungsfläche zwischen zweien solchen Räumen, in welchen, jeden für sich genommen, die Dichtigkeit nach der Stetigkeit, aber beim Übergange aus dem einen in den andern sprungsweise sich ändert, so haben daselbst, allgemein zu reden, $\frac{\mathrm{d}\mathrm{d}V}{\mathrm{d}x^2}$, $\frac{\mathrm{d}\mathrm{d}V}{\mathrm{d}y^2}$, $\frac{\mathrm{d}\mathrm{d}V}{\mathrm{d}z^2}$ jedes zwei verschiedene Werthe, und von dem Aggregate jener Grössen gilt dasselbe, was am Schlusse des 8. Artikels erinnert ist.
\SourceLeaf{218}%
\Article{12}
Wir ziehen, wie schon oben bemerkt ist, auch den idealen Fall mit in den Kreis unsrer Untersuchungen, wo Anziehungs- oder Abstossungskräfte von den Theilen einer \textit{Fläche} ausgehend angenommen werden, und erlauben uns dabei die Einkleidung, dass eine wirkende Masse in der Fläche vertheilt sei. Unter Dichtigkeit in irgend einem Punkte der Fläche verstehen wir in diesem Falle den Quotienten, wenn die in einem Elemente der Fläche, welchem der Punkt angehört, enthaltene Masse mit diesem Elemente dividirt wird. Diese Dichtigkeit kann gleichförmig (in allen Punkten dieselbe) oder ungleichförmig sein, und im letztern Falle entweder in der ganzen Fläche sich nach der Stetigkeit ändern (d. i. so, dass sie in je zwei einander unendlich nahen Punkten auch nur unendlich wenig verschieden ist) oder es kann die ganze Fläche in zwei oder mehrere Stücke zerfallen, in deren jedem eine stetige Änderung Statt findet, während beim Übergange aus einem in das andere die Änderung sprungsweise geschieht. Übrigens kann auch eine solche Vertheilung gedacht werden, wo unbeschadet der Endlichkeit der ganzen Masse, die Dichtigkeit in einzelnen Punkten oder Linien unendlich gross wird. Der Fläche selbst, insofern sie nicht eine Ebene ist, wird allgemein zu reden eine stetige Krümmung beigelegt werden, ohne darum eine Unterbrechung in einzelnen Punkten (Ecken) oder Linien (Kanten) auszuschliessen.

\expandafter\def\csname HistoricalNoteText22\endcsname{Von der Endlichkeit des Integrals, welches das Potential ausdrückt, überzeugt man sich leicht, indem man die Zerlegung der Fläche in Elemente auf ähnliche Weise ausführt, wie im 15. Artikel geschehen wird; und zugleich wird daraus ersichtlich, dass die den beiden in Rede stehenden Punkten unendlich nahen Theile der Fläche zu dem ganzen Integral nur unendlich wenig beitragen, woraus sich das oben gesagte leicht beweisen lässt.}%
Dies vorausgesetzt erhält das Potential auch in jedem Punkte der Fläche selbst, wo nur die Dichtigkeit nicht unendlich gross ist, einen bestimmten endlichen Werth, von welchem der Werth in einem zweiten Punkt, der, in der Fläche oder ausserhalb, jenem unendlich nahe liegt, nur unendlich wenig verschieden sein kann\HistoricalNote{22}, oder mit andern Worten, in jeder Linie, möge sie in der Fläche selbst liegen, oder dieselbe kreuzen, ändert sich das Potential nach der Stetigkeit.

\Article{13}
Bezeichnet man mit $k$ die Dichtigkeit in dem Flächenelement $\mathrm{d}s$; mit
\SourceLeaf{219}%
$a$, $b$, $c$ die Coordinaten eines demselben angehörenden Punkts; mit $r$ dessen Entfernung von einem Punkte $O$, dessen Coordinaten $x$, $y$, $z$ sind, und mit $V$ das Potential der in der Fläche enthaltenen Masse in dem Punkte $O$, so ist $V=\int\frac{k\mathrm{d}s}{r}$, durch die ganze Fläche ausgedehnt, endlich mit $X$, $Y$, $Z$ die eben so verstandenen Integrale
\[
\int\frac{k(a-x)\mathrm{d}s}{r^3},\quad\int\frac{k(b-y)\mathrm{d}s}{r^3},\quad\int\frac{k(c-z)\mathrm{d}s}{r^3}
\]
so sind zwar $X$, $Y$, $Z$ ganz gleichbedeutend mit $\frac{\mathrm{d}V}{\mathrm{d}x}$, $\frac{\mathrm{d}V}{\mathrm{d}y}$, $\frac{\mathrm{d}V}{\mathrm{d}z}$, solange $O$ ausserhalb der Fläche liegt, aber genau zu reden gilt dies nicht mehr, wenn $O$ ein Punkt der Fläche selbst ist, und die Ungleichheit gestaltet sich verschieden je nach der Beschaffenheit des Winkels, welchen die Normale auf die Fläche mit der betreffenden Coordinatenaxe macht. Es ist offenbar hinreichend, hier nur das Verhalten in Beziehung auf die erste Coordinatenaxe anzugeben.

I. Ist jener Winkel $=0$, so hat in $O$ das Integral $X$ einen bestimmten Werth, $\frac{\mathrm{d}V}{\mathrm{d}x}$ hingegen hat zwei verschiedene Werthe, je nachdem man $\mathrm{d}x$ als positiv oder als negativ betrachtet.

II. Ist der Winkel ein rechter, so lässt der Ausdruck für $X$ eine wahre Integration nicht zu (indem dann eine ähnliche Bemerkung gilt, wie im 7. Artikel), während $\frac{\mathrm{d}V}{\mathrm{d}x}$ nur Einen bestimmten Werth hat.

III. Ist der Winkel spitz, so verhält es sich mit $X$ eben so wie im zweiten, und mit $\frac{\mathrm{d}V}{\mathrm{d}x}$ eben so wie im ersten Falle.

Noch besondere Modificationen treten ein, wenn in $O$ eine Unterbrechung der Stetigkeit entweder in Beziehung auf die Dichtigkeit oder die Krümmung Statt findet. Für unsern Hauptzweck ist jedoch nicht nothwendig, solche Ausnahmsfälle, die nur in einzelnen Linien oder Punkten eintreten können, ausführlich abzuhandeln, und wir werden daher bei der nähern Erörterung des Gegenstandes annehmen, dass in dem fraglichen Punkte eine bestimmte endliche Dichtigkeit, und eine bestimmte Berührungsebene Statt findet.

\Article{14}
Ehe wir die Untersuchung in ihrer Allgemeinheit vornehmen, wird es nützlich sein, einen einfachen besondern Fall zu betrachten. Es sei die Fläche das Stück $A$ einer Kugelfläche, und die Dichtigkeit darin gleichförmig oder $k$ constant. Es sind also $V$, $X$ die Werthe der Integrale
\SourceLeaf{220}%
\[
\int\frac{k\mathrm{d}s}{r},\quad\int\frac{k(a-x)\mathrm{d}s}{r^3}
\]
durch $A$ ausgedehnt; bezeichnen wir mit $V'$, $X'$ dieselben Integrale, wenn sie durch den übrigen Theil der Kugelfläche $B$, und mit $V^0$, $X^0$, wenn sie durch die ganze Kugelfläche erstreckt werden, so wird $V=V^0-V'$, $X=X^0-X'$. Wir wollen noch den Halbmesser der Kugel mit $R$ bezeichnen, den Anfangspunkt der Coordinaten in den Mittelpunkt der Kugel legen, und $\BareRoot{(xx+yy+zz)}$ oder den Abstand des Punktes $O$ vom Mittelpunkte der Kugel $=\rho$ setzen.

Es ist nun bekannt, dass $V^0=4\pi kR$ wird, wenn $O$ innerhalb der Kugel, hingegen $V^0=\frac{4\pi kRR}{\rho}$, wenn $O$ ausserhalb liegt; in der Kugelfläche selbst fallen beide Werthe zusammen. Der Differentialquotient $\frac{\mathrm{d}V^0}{\mathrm{d}x}$ wird daher innerhalb der Kugel $=0$, ausserhalb $=-\frac{4\pi kRRx}{\rho^3}$; auf der Kugelfläche selbst aber werden beide Werthe zugleich gelten, je nach dem Zeichen von $\mathrm{d}x$: gleich sind diese beiden Werthe nur dann, wenn $x=0$ ist, was dem Falle II des vorhergehenden Artikels entspricht.

Der Ausdruck für $X^0$, innerhalb und ausserhalb der Kugel mit $\frac{\mathrm{d}V^0}{\mathrm{d}x}$ gleichbedeutend, wird auf der Oberfläche ein leeres Zeichen, insofern eine wahre Integration unstatthaft ist, den einzigen Fall ausgenommen, wenn für die unendlich nahe liegenden Elemente der Fläche $a-x$ ein unendlich kleines von einer höhern Ordnung wird als $r$, nemlich wenn $y=0$, $z=0$. $x=\pm R$, für welchen Fall die Integration $X^0=\mp2\pi k$ gibt, also mit keinem der Werthe von $\frac{\mathrm{d}V^0}{\mathrm{d}x}$ übereinstimmend, sondern vielmehr mit dem Mittel von beiden: offenbar gehört übrigens dieser Fall zu I im vorhergehenden Artikel.

Erwägt man nun, dass wenn $O$ ein auf der Oberfläche der Kugel innerhalb $A$ liegender Punkt ist, $X'$ und $\frac{\mathrm{d}V'}{\mathrm{d}x}$ gleichbedeutend sind und bestimmte nach der Stetigkeit sich ändernde Werthe haben, so erhellt, dass das gegenseitige Verhalten zwischen $X^0-X'$ und $\frac{\mathrm{d}V^0}{\mathrm{d}x}-\frac{\mathrm{d}V'}{\mathrm{d}x}$, d. i. zwischen $X$ und $\frac{\mathrm{d}V}{\mathrm{d}x}$ ganz dasselbe ist, wie zwischen $X^0$ und $\frac{\mathrm{d}V^0}{\mathrm{d}x}$, woraus also die im vorhergehenden Artikel aufgestellten Sätze von selbst folgen.

\Article{15}
Für die allgemeinere Untersuchung ist es vortheilhaft, den Anfangspunkt der Coordinaten in einen in der Fläche selbst liegenden Punkt $P$ zu setzen, und die erste Coordinatenaxe senkrecht gegen die Berührungsebene in $P$ zu legen.
\SourceLeaf{221}%

Bezeichnen wir mit $\psi$ den Winkel zwischen der Normale auf das unbestimmte Flächenelement $\mathrm{d}s$ und der ersten Coordinatenaxe, so ist $\cos\psi.\mathrm{d}s$ die Projection von $\mathrm{d}s$ auf die Ebene der $b$ und $c$; und setzen wir
\[
\BareRoot{(bb+cc)}=\rho,\quad b=\rho\cos\theta,\quad c=\rho\sin\theta
\]
so wird $\rho\mathrm{d}\rho.\mathrm{d}\theta$ ein unbestimmtes Element dieser Ebene vorstellen, und das entsprechende Flächenelement $\mathrm{d}s=\frac{\rho\mathrm{d}\rho.\mathrm{d}\theta}{\cos\psi}$ sein; das darin enthaltene Massenelement wird also $=h\rho\mathrm{d}\rho.\mathrm{d}\theta$ sein, wenn wir zur Abkürzung $h$ für $\frac{k}{\cos\psi}$ schreiben.

Wir wollen nun untersuchen, inwiefern der Werth von $X$ sich sprungsweise ändert, indem der Punkt $O$ in der ersten Coordinatenaxe von der einen Seite der Fläche auf die andere, oder $x$ aus einem negativen Werthe in einen positiven übergeht. Für diese Frage ist es offenbar einerlei, ob wir die ganze Fläche in Betracht ziehen, oder nur einen beliebig kleinen, den Punkt $P$ einschliessenden Theil, da der Beitrag des übrigen Theils der Fläche zu dem Werthe von $X$ sich nach der Stetigkeit ändert. Es ist daher erlaubt, $\rho$ nur von $0$ bis zu einem beliebig kleinen Grenzwerthe $\rho'$ auszudehnen, und vorauszusetzen, dass in der so begrenzten Fläche $h$ und $\frac{a}{\rho}$ sich überall nach der Stetigkeit ändern. Setzen wir, für jeden bestimmten Werth von $\theta$, den Werth des Integrals $\int\frac{h(a-x)\rho\mathrm{d}\rho}{r^3}$, von $\rho=0$ bis $\rho=\rho'$ ausgedehnt, $=Q$, so wird $X=\int Q\mathrm{d}\theta$, wo die Integration von $\theta=0$, bis $\theta=2\pi$ zu erstrecken ist.

Es kommt nun darauf an, die Werthe von $X$ für $x=0$, für ein unendlich kleines positives $x$, und für ein unendlich kleines negatives (die beiden andern Coordinaten $y$, $z$ allemal $=0$ angenommen) unter einander zu vergleichen; wir bezeichnen diese drei Werthe von $X$ mit $X^0$, $X'$, $X''$, und die entsprechenden Werthe von $Q$ mit $Q^0$, $Q'$, $Q''$.

Da $r=\BareRoot{((a-x)^2+\rho\rho)}$, so erhält man, indem man $\theta$ als constant betrachtet
\[
\mathrm{d}\frac{h(a-x)}{r}=-\frac{h(a-x)\rho\mathrm{d}\rho}{r^3}+\frac{\mathrm{d}h}{\mathrm{d}\rho}.\frac{a-x}{r}.\mathrm{d}\rho+\frac{\mathrm{d}a}{\mathrm{d}\rho}.\frac{h\rho\rho}{r^3}.\mathrm{d}\rho
\]
und folglich
\[
Q=\int\frac{\mathrm{d}h}{\mathrm{d}\rho}.\frac{a-x}{r}.\mathrm{d}\rho+\int\frac{\mathrm{d}a}{\mathrm{d}\rho}.\frac{h\rho\rho}{r^3}.\mathrm{d}\rho-\frac{h'(a'-x)}{r'}+\text{Const.}
\]
wo die beiden Integrationen von $\rho=0$ bis $\rho=\rho'$ auszudehnen, und die Werthe
